\documentclass[UTF8,a4paper,11pt]{ctexart}
\usepackage[a4paper,margin=2.3cm]{geometry}
\usepackage{amsmath,amssymb,mathtools}
\usepackage{enumitem}
\usepackage[hidelinks]{hyperref}
\usepackage{fancyhdr}
\setlist[itemize]{leftmargin=2.2em,itemsep=0.25em,topsep=0.35em}
\setlength{\parindent}{2em}\setlength{\parskip}{0.35em}
\pagestyle{fancy}\fancyhf{}\fancyhead[L]{离散数学与结构（I）}\fancyhead[R]{Lec 18}\fancyfoot[C]{\thepage}\setlength{\headheight}{14pt}
\title{互信息、集中不等式与 Sanov 大偏差}\author{zhangjunshuo}\date{2026 年 11 月}
\begin{document}
\maketitle
\begin{quote}
本讲根据 2026 年 11 月 17 日的三页手写稿、上半节旧转写和下半节重新转写整理。上一份笔记是 Lec17（11 月 12 日），本次按顺序记为 \textbf{Lec18}。手写稿是公式与结构的主要依据，转写用于补课堂解释与例子；下半节旧转写不完整，本笔记使用重新转写的版本。

本讲主线：KL/互信息 $\to$ Markov/Chebyshev/Chernoff $\to$ 类型计数与 Sanov 大偏差 $\to$ 数据处理不等式。

每个分节开头都先交代本节用到的记号；所有符号在首次使用处给出定义。
\end{quote}
\setcounter{tocdepth}{1}\tableofcontents\clearpage

\section{KL 散度、互信息与条件互信息}


这一讲从 KL 散度与熵的关系出发，引入\textbf{互信息}和\textbf{条件互信息}，再用它们统一描述“两个随机变量共享多少信息”。这部分是后面集中不等式、类型计数与 Sanov 大偏差的语言基础。

\textbf{本节用到的记号：} $H(X)$ 为熵；$H(X\mid Y)$ 为条件熵；$I(X;Y)$ 为互信息；$I(X;Y\mid Z)$ 为条件互信息；$D(P\Vert Q)$ 为 KL 散度；$U$ 为均匀分布。

\subsection{KL 散度与均匀分布}


设 $P$ 是 $[n]$ 上的分布，$U$ 是 $[n]$ 上的均匀分布。则

\[
\begin{aligned}
D(P\Vert U)
&=\sum_{x=1}^{n}P(x)\log\frac{P(x)}{1/n}\\
&=\sum_xP(x)\log n-\sum_xP(x)\log\frac1{P(x)}\\
&=\log n-H(P).
\end{aligned}
\]

也可以写成

\[
D(P\Vert U)=H(U)-H(P).
\]

\textbf{解释.} 均匀分布是“最大熵”分布；$D(P\Vert U)$ 正好衡量 $P$ 比均匀分布少了多少熵。若 $P=U$，散度为 $0$；若 $P$ 集中在少数点，散度变大。

\subsection{互信息的定义}


两个随机变量 $X,Y$ 的\textbf{互信息}定义为

\[
\boxed{
I(X;Y)
=H(X)+H(Y)-H(X,Y)
}.
\]

由链式法则 $H(X,Y)=H(X)+H(Y\mid X)$，得到几个等价形式：

\[
\begin{aligned}
I(X;Y)
&=H(X)-H(X\mid Y)\\
&=H(Y)-H(Y\mid X)\\
&=D(P_{X,Y}\Vert P_XP_Y).
\end{aligned}
\]

最后一条说明：互信息就是联合分布与“边缘分布乘积”之间的 KL 散度；它衡量 $X,Y$ 的“依赖程度”。

\textbf{基本性质.}

\begin{itemize}
\item 对称：$I(X;Y)=I(Y;X)$。
\item 非负：$I(X;Y)\ge0$；等号当且仅当 $X,Y$ 独立。
\item $I(X;X)=H(X)$。
\item $I(X;Y)\le H(X)$ 且 $I(X;Y)\le H(Y)$。
\end{itemize}


\textbf{Venn 图直观.} 把 $H(X)$ 和 $H(Y)$ 看成两个集合，交叠部分是 $I(X;Y)$，两边分别是 $H(X\mid Y)$ 和 $H(Y\mid X)$，并集是 $H(X,Y)$。

\subsection{条件互信息}


给定 $Z$ 后，$X,Y$ 的\textbf{条件互信息}定义为

\[
\boxed{
I(X;Y\mid Z)
=H(X\mid Z)+H(Y\mid Z)-H(X,Y\mid Z)
}.
\]

等价地，

\[
I(X;Y\mid Z)
=D(P_{X,Y\mid Z}\Vert P_{X\mid Z}P_{Y\mid Z}\mid P_Z),
\]

即“在每个 $Z=z$ 的条件下算互信息，再按 $P_Z(z)$ 加权平均”。

\textbf{链式法则.}

\[
I(X;Y,Z)=I(X;Z)+I(X;Y\mid Z).
\]

\textbf{条件互信息可以变大，也可以变小.} 这是它和普通互信息最容易混淆的地方。

\textbf{例 1（条件化后互信息变小）.} 令

\[
X=Y=Z\sim\mathrm{Bern}(1/2).
\]

则 $X$ 与 $Y$ 完全相关，$I(X;Y)=1$ bit；但给定 $Z$ 后 $X,Y$ 都等于 $Z$，没有剩余不确定性，所以

\[
I(X;Y\mid Z)=0.
\]

\textbf{例 2（条件化后互信息变大）.} 令 $X,Y$ 独立同分布 $\mathrm{Bern}(1/2)$，令

\[
Z=X\oplus Y.
\]

则 $X,Y$ 独立，$I(X;Y)=0$；但给定 $Z$ 后，知道 $X$ 就能确定 $Y$，所以

\[
I(X;Y\mid Z)=1\text{ bit}.
\]

\textbf{例 3（explaining away）.} 若 $Z$ 是两个原因 $X,Y$ 共同影响的观察结果，则 $X,Y$ 可能在边际上独立，但给定 $Z$ 后变得相关：一个原因解释了观察结果，另一个原因的可能性就下降。这正是“条件互信息大于边际互信息”的典型场景。

\subsection{互信息与熵的关系}


常用的等价关系：

\[
H(X,Y)=H(X)+H(Y)-I(X;Y),
\]

\[
H(X\mid Y)=H(X)-I(X;Y).
\]

因此：

\begin{itemize}
\item $I(X;Y)=0$ $\iff$ $X,Y$ 独立；
\item $I(X;Y)=H(X)$ $\iff$ $H(X\mid Y)=0$，即给定 $Y$ 后 $X$ 没有不确定性（$X$ 是 $Y$ 的函数）。
\end{itemize}


\subsection{常见错误}


\begin{itemize}
\item 把 $I(X;Y\mid Z)$ 直接当成 $I(X;Y)$；条件化会改变依赖结构，可能增大也可能减小。
\item 忘记互信息非负；算出负数说明公式或分布写错了。
\item 把 $I(X;Y)=0$ 当成“$X,Y$ 完全不相关”的弱结论；它实际上等价于 $X,Y$ 独立，条件很强。
\item 把 $H(X\mid Y)$ 与 $H(X,Y)$ 混淆：前者是“已知 $Y$ 后 $X$ 的剩余熵”，后者是联合熵。
\end{itemize}

\section{集中不等式：Markov、Chebyshev 与 Chernoff}


这一节把“偏离期望的概率有多小”变成可计算的不等式。三个层次依次是：

1. \textbf{Markov}：只用非负性；
2. \textbf{Chebyshev}：用方差；
3. \textbf{Chernoff}：用矩生成函数/指数变换，得到指数级尾界。

\textbf{本节用到的记号：} $X$ 为随机变量；$\mu=E[X]$；$\Delta=\sqrt{\operatorname{Var}(X)}$；$D(q\Vert p)$ 为二元 KL；$X_1,\dots,X_n$ 独立同分布。

\subsection{Markov 不等式}


若 $X\ge0$，则对任意 $a>0$，

\[
\boxed{
P(X\ge a)\le\frac{E[X]}{a}
}.
\]

等价地，取 $a=\alpha E[X]$：

\[
P(X\ge\alpha E[X])\le\frac1\alpha\qquad(\alpha>0).
\]

\textbf{证明.} 因为 $X\ge a\mathbf 1_{\{X\ge a\}}$，取期望：

\[
E[X]\ge aP(X\ge a).
\]

两边除以 $a$ 即得。

\textbf{用途.} Markov 只要求非负，条件最弱，但界通常很松；它是后面 Chebyshev 和 Chernoff 的证明工具。

\subsection{Chebyshev 不等式}


对任意 $a>0$，

\[
\boxed{
P(|X-E[X]|\ge a)\le\frac{\operatorname{Var}(X)}{a^2}
}.
\]

若记 $\Delta=\sqrt{\operatorname{Var}(X)}$，则

\[
P(|X-E[X]|\ge a\Delta)\le\frac1{a^2}.
\]

\textbf{证明.} 对非负随机变量 $(X-E[X])^2$ 用 Markov：

\[
P((X-E[X])^2\ge a^2)
\le\frac{E[(X-E[X])^2]}{a^2}
=\frac{\operatorname{Var}(X)}{a^2}.
\]

\textbf{用途.} 只要知道方差，就能给出“偏离 $a$ 个标准差”的概率上界；但尾界是多项式级（$1/a^2$），不够精细。

\subsection{Chernoff 不等式}


对任意 $a$ 和任意 $t>0$，

\[
\boxed{
P(X\ge a)
\le\min_{t>0}\frac{E[e^{tX}]}{e^{ta}}
}.
\]

\textbf{证明.} 对非负随机变量 $e^{tX}$ 用 Markov：

\[
P(X\ge a)=P(e^{tX}\ge e^{ta})
\le\frac{E[e^{tX}]}{e^{ta}}.
\]

对所有 $t>0$ 取最小，得到最紧的界。

\textbf{为什么强.} $E[e^{tX}]$ 把整个分布的信息都编码进来；选择最优 $t$ 后，尾界通常是指数级 $e^{-\Theta(n)}$，比 Chebyshev 的 $1/a^2$ 精细得多。

\subsection{Bernoulli 和的 Chernoff 界}


设 $X_1,\dots,X_n$ 独立同分布 $\mathrm{Bern}(p)$，$q>p$。则

\[
P\Bigl(\sum_{i=1}^nX_i\ge qn\Bigr)
\le\exp\bigl(-nD(q\Vert p)\bigr),
\]

其中

\[
D(q\Vert p)
=q\log\frac qp+(1-q)\log\frac{1-q}{1-p}.
\]

\textbf{推导.} 对任意 $t>0$，

\[
\begin{aligned}
P\Bigl(\sum_iX_i\ge qn\Bigr)
&\le\frac{E[e^{t\sum_iX_i}]}{e^{tqn}}\\
&=\frac{\bigl(pe^t+(1-p)\bigr)^n}{e^{tqn}}.
\end{aligned}
\]

令 $f(t)=\log\bigl(pe^t+(1-p)\bigr)-tq$，求导并令其为零：

\[
\frac{pe^t}{pe^t+(1-p)}=q
\quad\Longrightarrow\quad
e^t=\frac{q(1-p)}{p(1-q)}.
\]

代回并整理，指数变为

\[
-\Bigl[q\log\frac qp+(1-q)\log\frac{1-q}{1-p}\Bigr]
=-D(q\Vert p).
\]

\textbf{二次下界.} 采用自然对数时，当 $q,p$ 接近，

\[
D(q\Vert p)\ge 2(p-q)^2.
\]

因此

\[
P\Bigl(\sum_iX_i\ge qn\Bigr)
\le\exp\bigl(-2n(p-q)^2\bigr).
\]

若改用 $\log_2$，$D$ 的数值会放大 $1/\ln2$ 倍，$\exp$ 也要换成以 $2$ 为底；不要混用底数。手写稿里的 $2\log e$ 就是换底因子。

\subsection{与 CLT / 正态尾的比较}


对 $p=1/2$，中心极限定理给出

\[
\frac{\sum_iX_i-n/2}{\sqrt{n}/2}
\approx \mathcal N(0,1).
\]

因此 $P(\sum_iX_i\ge qn)$ 的正态近似尾概率约为

\[
\exp\bigl(-2n(q-1/2)^2\bigr),
\]

与 Chernoff 界的指数一致。

\textbf{但要区分“近似”和“上界”：}

\begin{itemize}
\item CLT 给的是\textbf{近似}，不是严格的尾上界；
\item 严格的指数上界用 Chernoff；
\item 若要一个简单的严格界，也可用 Hoeffding：
\end{itemize}

\[
  P(|\bar X-E[X]|\ge a)\le2e^{-2na^2}.
\]

手写稿把“CLT 正态尾”和“Chernoff 界”放在一起比较，是为了说明两者的指数相同；证明严格上界时不要把 CLT 当成不等式使用。

\subsection{三种界的关系}


\begin{itemize}
\item \textbf{Markov}：需要“非负、期望”；尾界量级 $1/a$；最粗略。
\item \textbf{Chebyshev}：需要“方差”；尾界量级 $1/a^2$；快速估计。
\item \textbf{Chernoff}：需要“矩生成函数”；尾界量级 $e^{-\Theta(n)}$；独立和、大偏差。
\end{itemize}


\textbf{记忆.} Markov 看“质量”，Chebyshev 看“方差”，Chernoff 看“指数矩”；条件越强，界越紧。

\subsection{常见错误}


\begin{itemize}
\item 对非正随机变量直接用 Markov。
\item 把 Chernoff 写成 $P(X\ge a)\le e^{-ta}E[e^{tX}]$ 后忘记对所有 $t>0$ 取最小。
\item 把 CLT 的近似尾概率当成严格上界。
\item 在 Bernoulli 界中把 $D(q\Vert p)$ 的两个参数顺序写反：$D(q\Vert p)$ 是“用 $p$ 近似 $q$”的代价，不是 $D(p\Vert q)$。
\end{itemize}

\section{类型计数与 Sanov 大偏差}


这一节把“大偏差事件的指数”从 Bernoulli 推广到一般有限字母表：先数“类型”（经验分布），再用类型类的大小和概率得到 \textbf{Sanov 界}。

\textbf{本节用到的记号：} 字母表 $[k]=\{1,\dots,k\}$；$P$ 为真实分布；$X_1,\dots,X_n$ iid $\sim P$；$T_i=\#\{j:X_j=i\}$ 为频数；$Q_i=T_i/n$ 为经验分布/类型；$H(Q)$ 为熵；$D(Q\Vert P)$ 为 KL 散度。

\subsection{Bernoulli 尾概率的精确写法}


先看 $k=2$ 的情形：$X_i\sim\mathrm{Bern}(p)$。对 $q>p$，

\[
P\Bigl(\sum_{i=1}^nX_i\ge qn\Bigr)
=\sum_{m=\lceil qn\rceil}^{n}\binom nm p^m(1-p)^{n-m}.
\]

\textbf{上界：用最大项.} 和式不超过“项数 × 最大项”：

\[
P\Bigl(\sum_iX_i\ge qn\Bigr)
\le n\max_{m\ge qn}\binom nm p^m(1-p)^{n-m}.
\]

在 $m\ge qn>pn$ 的区域，二项概率的最大值在 $m=\lceil qn\rceil$ 附近；用 Stirling 公式

\[
n!\approx\sqrt{2\pi n}\left(\frac ne\right)^n
\]

可得

\[
\binom nm\approx\exp\bigl(nH(m/n)\bigr),
\]

其中 $H(x)=-x\log x-(1-x)\log(1-x)$。于是最大项约为

\[
\exp\bigl(nH(q)\bigr)p^{qn}(1-p)^{(1-q)n}
=\exp\bigl(-nD(q\Vert p)\bigr).
\]

因此

\[
P\Bigl(\sum_iX_i\ge qn\Bigr)
\le n\exp\bigl(-nD(q\Vert p)\bigr).
\]

\textbf{下界：只留 $m=qn$ 一项.}

\[
P\Bigl(\sum_iX_i\ge qn\Bigr)
\ge\binom{n}{qn}p^{qn}(1-p)^{(1-q)n}.
\]

这是 $k=2$ 时类型类大小的标准下界：类型类 $T(Q)$ 的大小不少于

\[
\frac{\exp(nH(Q))}{n+1}.
\]

取 $Q=(q,1-q)$，即得

\[
\binom{n}{qn}\ge\frac1{n+1}\exp\bigl(nH(q)\bigr).
\]

于是

\[
P\Bigl(\sum_iX_i\ge qn\Bigr)
\ge\frac1{n+1}\exp\bigl(-nD(q\Vert p)\bigr).
\]

\textbf{结论（Bernoulli 版 Sanov）.}

\[
\boxed{
\frac1{n+1}e^{-nD(q\Vert p)}
\le
P\Bigl(\sum_iX_i\ge qn\Bigr)
\le
n\,e^{-nD(q\Vert p)}
}.
\]

上下界只差多项式因子，\textbf{指数 $D(q\Vert p)$ 是紧的}。

\subsection{类型与类型类}


一般字母表 $[k]$。给定一个序列 $x=(x_1,\dots,x_n)$，它的\textbf{类型}是经验分布

\[
Q_i=\frac{T_i}{n},\qquad T_i=\#\{j:x_j=i\}.
\]

所有具有同一类型 $Q$ 的序列组成\textbf{类型类}，记作 $T(Q)$。

\textbf{类型类大小.}

\[
|T(Q)|=\frac{n!}{T_1!\cdots T_k!}.
\]

用 Stirling 公式可得

\[
|T(Q)|\approx\exp\bigl(nH(Q)\bigr),
\]

更精确地有

\[
\exp\bigl(nH(Q)\bigr)\big/(n+1)^{k-1}
\le |T(Q)|
\le \exp\bigl(nH(Q)\bigr).
\]

\textbf{单个序列的概率.}

\[
P^n(x)=\prod_{i=1}^kP_i^{T_i}
=\exp\Bigl(n\sum_iQ_i\log P_i\Bigr).
\]

因此类型类 $T(Q)$ 的总概率为

\[
P^n(T(Q))
=|T(Q)|\prod_iP_i^{nQ_i}
\approx\exp\bigl(-nD(Q\Vert P)\bigr),
\]

其中

\[
D(Q\Vert P)=\sum_{i=1}^kQ_i\log\frac{Q_i}{P_i}.
\]

\textbf{精确的上下界}（差多项式因子）：

\[
\frac{1}{(n+1)^{k-1}}e^{-nD(Q\Vert P)}
\le P^n(T(Q))
\le e^{-nD(Q\Vert P)}.
\]

\subsection{Sanov 界}


设 $\mathcal R$ 是“类型/经验分布”的一个集合（例如 $\{Q:Q_1\ge q\}$）。则

\[
\boxed{
P^n(Q\in\mathcal R)
\le (n+1)^{k-1}
\exp\Bigl(-n\min_{Q\in\mathcal R}D(Q\Vert P)\Bigr)
}.
\]

并有下界

\[
P^n(Q\in\mathcal R)
\ge \frac{1}{(n+1)^{k-1}}
\exp\Bigl(-n\min_{Q\in\mathcal R}D(Q\Vert P)\Bigr)
\]

（在 $\mathcal R$ 中取到最小值附近的类型时）。

令

\[
Q^*=\arg\min_{Q\in\mathcal R}D(Q\Vert P),
\]

则

\[
P^n(Q\in\mathcal R)
=\exp\Bigl(-nD(Q^*\Vert P)+o(n)\Bigr).
\]

\textbf{直观.} 大偏差事件的指数，由“区域 $\mathcal R$ 中离真实分布 $P$ 最近的那个分布 $Q^*$”决定；$D(Q^*\Vert P)$ 就是偏离的“代价”。

\subsection{Bernoulli 情形是特例}


取 $k=2$、$\mathcal R=\{Q:Q_1\ge q\}$，则

\[
Q^*=(q,1-q),\qquad
D(Q^*\Vert P)=D(q\Vert p).
\]

于是 Sanov 界退化为前面的

\[
P\Bigl(\sum_iX_i\ge qn\Bigr)
\asymp\exp\bigl(-nD(q\Vert p)\bigr),
\]

与 Chernoff 的指数一致；Sanov 还给出下界，因此说明 Chernoff 的指数是紧的。

\subsection{一般区域的几何图像}


把分布空间画成一个区域，$P$ 是真实分布，$\mathcal R$ 是关心的“坏事件”区域：

\begin{itemize}
\item 若 $\mathcal R$ 不包含 $P$，大偏差概率指数衰减；
\item 衰减速率 = 从 $P$ 到 $\mathcal R$ 的“KL 距离”；
\item $Q^*$ 是 $\mathcal R$ 中离 $P$ 最近的点；在连续/正则区域中，它通常位于区域边界上。
\end{itemize}


\textbf{注意.} 若 $\mathcal R$ 取得很“碎”（例如只包含无理频率点），需要额外说明最小值的取法；课堂上加了“频率可实现”的条件，避免这种退化情形。

\subsection{与 Chernoff 的分工}


\begin{itemize}
\item \textbf{Chernoff}：给上界，计算通常更直接，适合“独立和超过某阈值”的尾概率；
\item \textbf{Sanov}：给上界 + 下界，说明指数紧，适合一般的大偏差区域；
\item 在 Bernoulli 和一般有限字母表上，两者的指数都是 KL 散度 $D(Q^*\Vert P)$。
\end{itemize}


\subsection{常见错误}


\begin{itemize}
\item 把 Sanov 的 $Q^*$ 写成 $P$；$Q^*$ 是区域 $\mathcal R$ 中离 $P$ 最近的分布，不一定等于 $P$。
\item 忽略多项式因子：Sanov 的指数精确，但上下界差 $(n+1)^{k-1}$ 级别的因子。
\item 把 $H(Q)$ 与 $H(P)$ 混淆：类型类大小用 $H(Q)$，概率指数用 $D(Q\Vert P)$。
\item 把 Bernoulli 的 $D(q\Vert p)$ 与 $D(p\Vert q)$ 写反。
\end{itemize}

\section{数据处理不等式（DPI）}


数据处理不等式回答一个很自然的问题：\textbf{对数据再做处理，会不会增加其中关于原始变量的信息？} 答案是不会。信息只会丢失或保持不变，不会凭空增加。

\textbf{本节用到的记号：} 马尔可夫链 $X\to Y\to Z$；$I(X;Y)$ 为互信息；$D(P\Vert Q)$ 为 KL 散度；$\Delta_{\mathrm{TV}}$ 为 TV 距离；信道 $P_{Y\mid X}$。

\subsection{马尔可夫链与条件独立}


称 $X\to Y\to Z$ 构成\textbf{马尔可夫链}，如果

\[
P_{XYZ}(x,y,z)
=P_X(x)P_{Y\mid X}(y\mid x)P_{Z\mid Y}(z\mid y).
\]

等价地：

\begin{quote}
给定 $Y$ 后，$X$ 与 $Z$ 条件独立。
\end{quote}


也记作 $X\perp Z\mid Y$。

\textbf{直观.} $X$ 是原始信息，$Y$ 是经过一次处理后的数据，$Z$ 是在 $Y$ 上再处理得到的数据；$Z$ 只能通过 $Y$ 依赖 $X$。

\subsection{互信息的数据处理不等式}


若 $X\to Y\to Z$，则

\[
\boxed{
I(X;Z)\le I(X;Y)
}.
\]

\textbf{证明（链式法则 + 条件独立）.} 由互信息的链式法则：

\[
I(X;Y,Z)=I(X;Y)+I(X;Z\mid Y).
\]

另一方面，

\[
I(X;Y,Z)=I(X;Z)+I(X;Y\mid Z).
\]

因为 $X\perp Z\mid Y$，所以

\[
I(X;Z\mid Y)=0.
\]

于是

\[
I(X;Y)=I(X;Z)+I(X;Y\mid Z)\ge I(X;Z).
\]

最后一步用了条件互信息非负。

\textbf{直观.} $Y$ 是 $X$ 的“一手信息”，$Z$ 是从 $Y$ 加工出来的“二手信息”；二手信息关于 $X$ 的含量不可能超过一手信息。

\subsection{KL 的数据处理不等式}


若 $Y,Z$ 是由同一个信道 $P_{Y\mid X}$ 分别从 $X$ 和 $X'$ 得到的，则

\[
D(P_Y\Vert Q_Y)\le D(P_X\Vert Q_X).
\]

也就是说，经过同一个信道处理后，两个分布只会更接近，不会更远。

\subsection{TV 距离的数据处理不等式}


TV 距离也有同样的收缩性质：

\[
\boxed{
\Delta_{\mathrm{TV}}(P_Y,Q_Y)
\le
\Delta_{\mathrm{TV}}(P_X,Q_X)
}.
\]

\textbf{证明（耦合）.} 取 $P_X,Q_X$ 的最优耦合 $(X,X')$，使得

\[
\Pr[X\ne X']=\Delta_{\mathrm{TV}}(P_X,Q_X).
\]

对 $X,X'$ 分别独立地通过同一个信道 $P_{Y\mid X}$，得到 $Y,Y'$。则 $Y$ 的边缘是 $P_Y$，$Y'$ 的边缘是 $Q_Y$，且

\[
\Pr[Y\ne Y']\le\Pr[X\ne X'].
\]

而 TV 距离是“任意耦合下不相等的概率”的最小值，所以

\[
\Delta_{\mathrm{TV}}(P_Y,Q_Y)
\le\Pr[Y\ne Y']
\le\Delta_{\mathrm{TV}}(P_X,Q_X).
\]

\textbf{直观.} 同一个噪声信道只会把两个分布“搅得更像”；处理越多，越难区分。

\subsection{特例：确定性处理}


若 $Z=f(Y)$ 是确定性函数，则 $X\to Y\to Z$ 自动成立，于是

\[
I(X;f(Y))\le I(X;Y).
\]

这也解释了“对数据做任何确定性变换，都不会增加关于 $X$ 的信息”。

\subsection{常见错误}


\begin{itemize}
\item 把马尔可夫链方向写反：$X\to Y\to Z$ 表示给定 $Y$ 后 $X\perp Z$，不是给定 $Z$ 后 $X\perp Y$。
\item 认为 $I(X;Z)\le I(X;Y)$ 对任意联合分布都成立；它需要马尔可夫链/条件独立结构。
\item 把 TV 的 DPI 方向写反：处理之后 TV 只会变小或不变。
\item 认为“信息不会增加”意味着“信息一定严格减少”；确定性可逆处理可以保持不变。
\end{itemize}


\section*{一页式总览}
\begin{itemize}
\item KL 对均匀分布：$D(P\Vert U)=\log n-H(P)$。
\item 互信息：$I(X;Y)=H(X)+H(Y)-H(X,Y)=D(P_{XY}\Vert P_XP_Y)$。
\item 条件互信息：$I(X;Y\mid Z)=H(X\mid Z)+H(Y\mid Z)-H(X,Y\mid Z)$。
\item Markov：$P(X\ge a)\le E[X]/a$。
\item Chebyshev：$P(|X-E[X]|\ge a)\le\operatorname{Var}(X)/a^2$。
\item Chernoff：$P(X\ge a)\le\min_{t>0}E[e^{tX}]/e^{ta}$。
\item Bernoulli 和：$P(\sum_iX_i\ge qn)\le e^{-nD(q\Vert p)}$。
\item 类型：$Q_i=T_i/n$，$|T(Q)|\approx e^{nH(Q)}$，$P^n(T(Q))\approx e^{-nD(Q\Vert P)}$。
\item Sanov：$P(Q\in\mathcal R)\asymp e^{-n\min_{Q\in\mathcal R}D(Q\Vert P)}$。
\item DPI：$X\to Y\to Z$ $\Rightarrow$ $I(X;Z)\le I(X;Y)$。
\end{itemize}

\section*{参考资料}
\begin{itemize}
\item 课堂手写稿（2026-11-17，三页）。
\item 上半节旧转写 + 下半节重新转写（\texttt{课堂手记和录音转写/2026-11-17/}）。
\item \textbf{[DS]} 邓小铁、王畅。《离散数学与结构》信息论、集中不等式、Sanov 与数据处理不等式相关章节。
\end{itemize}
\end{document}
